%! Author = Saeid Yazdani
%! Date = 8/29/2024


\section{The FPGA Code for Hardware}
\label{sec:fpga-software}

The FPGA software is dedicated to the processing of real-time tasks.
It reads data from the measurement paths,
executes the control algorithms, and drives \acp{DAC} and \acp{ADC},
ensuring that the power supply maintains its desired output with precision.

\subsection{Control Register Module}
\label{subsec:fpga-control-registers}

As stated in \autoref{sec:fpga-mcu-structures}, the FPGA software interacts
with external and internal entities using registers.
That is, every setting, measurement data, output data, etc. has a single or
set of registers dedicated to it where other entities can read from or write to
these registers.

\autoref{lst:fpga-control-register-constants} shows how the register names,
addresses and reset values are encapsulated in a package for reusability and
sharing between other VHDL modules.

\begin{lstlisting}[language=VHDL, caption={Control registers global constants},
    label={lst:fpga-control-register-constants}]
package ctrl_register_pkg is

    -- Avalon BUS Attributes
    constant ADDR_WIDTH : natural :=  7;
    constant DATA_WIDTH : natural :=  16;
    -- ...

    -- Custom types
    type t_comp_array   is array ( natural range <> ) of std_logic_vector(15 downto 0);
    type t_timing_array is array ( natural range <> ) of std_logic_vector(c_tim_width-1 downto 0);
    -- ...

    -- Registers Indices
    constant V_ADC_IN    : natural := 0;
    constant V_ADC_SCALE : natural := 1;
    -- ...
    constant TEST_CONFIG : natural := 126;
    constant REVISION    : natural := 127;

    -- Registers Addresses
    constant ADDR_V_ADC_IN : std_logic_vector(ADDR_WIDTH-1 downto 0) := std_logic_vector(to_unsigned( V_ADC_IN    , ADDR_WIDTH ));
    -- ..

    -- Registers Reset Values
    constant ResetValue_V_ADC_IN    : std_logic_vector(DATA_WIDTH-1 downto 0) := x"0000";
    constant ResetValue_V_ADC_SCALE : std_logic_vector(DATA_WIDTH-1 downto 0) := x"1000";
    -- ...

end package;
\end{lstlisting}

Each register has an address and a reset value that will be applied on startup
to ensure system safety by disabling any output from the power supply.
\hyperref[sec:hvs-fpga-register-map]{Appendix A} contains the complete list
of registers used in the system.


The control register VHDL module manages various data transfer operations and
state transitions in the system, primarily by controlling read and write access
to multiple \acs{ADC} and \acs{DAC} registers which are essential for the
control loop operation.

As shown in \autoref{lst:fpga-control-register-entity}, the module is driven
by clock and reset signals and defines the inputs and outputs for the
communication with the microcontroller according to the \emph{AVALON} interface
followed by all the register port signals and their associated modes.

\begin{lstlisting}[language=VHDL, caption={Control registers entity},
    label={lst:fpga-control-register-entity}]
use ctrl_register_pkg.all;

entity ctrl_register is
    generic(
        g_aw            : integer := ADDR_WIDTH;
        g_dw            : integer := DATA_WIDTH;
        g_revision      : integer := 0
    );
    port (
        clk                   : in  std_logic;
        rst_n                 : in  std_logic;
        -- AVALON Interface for communication with the MCU
        chipselect      : IN  std_logic;
        address         : IN  std_logic_vector(g_aw -1 downto 0);
        write_en        : IN  std_logic;
        read_en         : IN  std_logic;
        writedata       : IN  std_logic_vector(g_dw -1 downto 0);
        readdata        : out std_logic_vector(g_dw -1 downto 0);
        waitrequest     : out std_logic;
        -- Registers
        V_ADC_DATA      : IN  std_logic_vector(g_dw -1 downto 0);
        V_ADC_DATA_EN   : IN  std_logic;
        V_ADC_IN        : OUT std_logic_vector(g_dw -1 downto 0);
        -- ...
        PA_COMP_ARRAY   : OUT t_comp_array  (0 to c_comp_width-1 );
        TIMING_ARRAY    : OUT t_timing_array(0 to c_comp_width-1 );
        -- Special output for test generation
        test_state      : out std_logic_vector(7 downto 0)
    );
end entity;
\end{lstlisting}

\autoref{lst:fpga-control-register-architecture} describes the \ac{RTL}
description of the control register entity where system registers are
updated at each clock.
When the system is powered up or on reset events, safe values of
the registers will be assigned.

The \emph{AVALON} state machine with various states such as \emph{av\_idle},
\emph{av\_read} and \emph{av\_write} is part of this \acs{RTL} description.
In the \emph{av\_idle} state, the module waits for a chip select signal to
initiate either a read or write operation, depending on the incoming control
signals.
When in the \emph{av\_read} state, the module reads data from specified
addresses corresponding to various ADC and DAC data registers and outputs it
through a data bus.
Similarly, in the \emph{av\_write} state, it writes incoming data to the
appropriate registers.

\pagebreak
\begin{lstlisting}[language=VHDL, caption={Control registers architecture},
    label={lst:fpga-control-register-architecture}]
architecture rtl of ctrl_register is

    -- Signals
    signal s_v_adc_in    : std_logic_vector(g_dw -1 downto 0);
    signal s_v_adc_scale : std_logic_vector(g_dw -1 downto 0);
    -- ...

begin

    -- Increment internal timers
    s_timer_high <= s_timer(g_dw*2-1 downto g_dw);
    s_timer_low  <= s_timer(g_dw-1 downto 0);

    p_main : process (clk,rst_n)

        -- AVALON Bus states
        type av_states is ( av_idle, av_rd_ready, av_wr_ready, av_rd_stop, av_wr_stop, av_write, av_read);

        variable av_state       : av_states;
        variable v_av_cnt       : unsigned(2 downto 0);

    begin

        if rst_n = '0' then -- Assign RESET values
            s_v_adc_in                  <=  ResetValue_V_ADC_IN;
            s_v_adc_scale               <=  ResetValue_V_ADC_SCALE;
            -- Reset other registers and states...

        elsif rising_edge(clk) then -- Perform clock based operations
        -- ...

            case av_state is
                when av_idle => -- Perform idle tasks
                when av_read => -- Perform read operation
                -- ...
            end case;
        end if;
    end process;

    -- Set registers
    V_ADC_IN                       <= s_v_adc_in;
    V_ADC_SCALE                    <= s_v_adc_scale;
    -- ...

end architecture;
\end{lstlisting}

The module also handles special control signals for enabling specific DAC
outputs and monitoring ADC values.
This design facilitates the interaction between the system's digital control
logic and its analog measurement components, enabling precise data acquisition
and output control in the overall system.

\subsection{ADC and DAC Modules}
\label{subsec:adc-dac-vhdl-modules}

The \ac{FPGA} issues clocks to \ac{ADC} chips on their \emph{CNV} pins to
begin conversion and then reads the resulting conversion to be used in
processing of the next system's update event.
The interfacing between \acs{ADC} chips and the \ac{FPGA} is done by a
dedicated \acs{VHDL} module as presented in \autoref{lst:fpga-ltc2386-entity}

\begin{lstlisting}[language=VHDL, caption={LTC2386 entity},
    label={lst:fpga-ltc2386-entity}]
ENTITY ltc2386 IS
    PORT(
        sys_clk    : in  std_logic;
        rst_n      : IN  std_logic;
        -- FPGA controller interface
        DATA       : OUT std_logic_vector(15 DOWNTO 0);
        DATA_EN    : OUT std_logic;
        -- LVDS inputs
        DCO        : IN  std_logic;
        DA         : IN  std_logic;
        DB         : IN  std_logic;
        -- LVDS outputs
        CLK_ADC    : OUT std_logic;
        CNV_EN     : OUT std_logic;
        -- TEST
        clk_smpl   : OUT std_logic
    );
END ENTITY;

ARCHITECTURE rtl OF ltc2386 IS

    SIGNAL dout         : unsigned(15 DOWNTO 0);
    -- ...
    SIGNAL s_cnv_en     : std_logic;

BEGIN

    DATA     <= std_logic_vector(dout) when rst_n = '1' else (others =>'0');
    data_en  <= '0' when rst_n = '0' else s_data_en ;
    -- ...

END ARCHITECTURE;
\end{lstlisting}

Similar to the \acs{ADC}, the \acs{DAC} output is controlled by the FPGA
and updates are issued by clocking the \emph{CLK} pins of \acs{DAC} chips.

\subsection{Compensation Module}
\label{subsec:compensation-vhdl-module}

The system is designed to control compensation network parameters based on
predefined arrays of compensation values and timing intervals.
The \emph{compensation\_ctrl} entity, defines the interface for the module,
including input signals for clock, reset, trigger, and arrays containing
the compensation values and timing data.

The output signals represent the current compensation settings for capacitors
and resistors in both the pre and power amplifiers
(two instances of the module).

Internally, the \acs{FSM} operates through a series of states: \emph{idle},
\emph{run}, \emph{set}, \emph{inc\_vi} and \emph{check\_vi}.
Upon reset, the system initializes the compensation settings to predefined
values and clears all counters and indexes.

The \acs{FSM} begins in the idle state, waiting for a rising edge on the
clock to start the process.
When triggered, the \acs{FSM} moves into the run state, where it checks if a
new rising edge has occurred and starts the process.
Otherwise, it waits until a specified timing interval has passed, as defined by
the timing array, before updating the compensation values in the set state.

In the set state, the \acs{FSM} loads the next set of compensation values
from the provided arrays and then increments the index to move to the next set
of values. It then checks if the index has reached the end of the arrays,
indicating that all compensation values have been applied.
If the end is reached, the \acs{FSM} returns to the idle state; otherwise, it
loops back to the run state to continue applying the next set of compensation
values.

The final compensation values are sent out through \acs{GPIO} pins of the
FPGA to the analog relays which control each resistor and capacitor in the
compensation network.

This system enables precise control over the application of compensation
parameters, allowing for fine-tuned adjustments in real-time based on the
predefined sequences and timing intervals based on \acs{DUT} requirements.

\autoref{lst:fpga-compensation-entity} in
\hyperref[sec:hvs-fpga-compensation-vhdl]{Appendix D} lists the complete VHDL
code responsible for driving the compensation networks.

\subsection{Control Loop Modules}
\label{subsec:fpga-control-loop-modules}

As described in \autoref{sec:loopdetails} and \autoref{fig:loop-algo},
multiple VHDL modules are working in a chain from reading the actual output
to generating the final output to the \emph{coarse}, \emph{vernier} and
\emph{set-ponit} DACs.

The averaging module performs a basic summing and division on ADC data
according to sample numbers set in the configuration register.
If the averaging is disabled or the sample size is set to one, the module
simply outputs the actual per-clock readings.

The calibration module will correct ADC's non-linearity and offset according
to the calibration table fed into the system by operator software
either at startup or during runtime.

After this stage scaling will be applied and the final output will be
sent out the output preparation module where the deviation from the set point
will be calculated and sent out on the appropriate DAC chip accordingly.

The code listings for the modules involved in control loop operations are
listed in \autoref{lst:fpga-calibration-entity},
\autoref{lst:fpga-scaling-entity} and \autoref{lst:fpga-outputprep-entity}
in \hyperref[sec:hvs-fpga-compensation-vhdl]{Appendix D}.









